\documentclass[11pt,a4paper]{article}
\usepackage[utf8]{inputenc}
\usepackage[T1]{fontenc}
\usepackage[english]{babel}
\usepackage{lmodern,microtype,graphicx,amsmath,geometry,xcolor,booktabs}
\usepackage[font=footnotesize,skip=3pt,labelsep=period]{caption}
\usepackage{titlesec}
\titleformat{\section}{\large\bfseries}{\thesection}{0.8em}{}
\titlespacing*{\section}{0pt}{6pt plus 1pt minus 1pt}{3pt}
\usepackage{hyperref}
\geometry{margin=2.0cm}
\hypersetup{colorlinks=true,linkcolor=blue!35!black,citecolor=blue,urlcolor=blue!35!black,pdftitle={Evaluating a Retraining Program When Its Losers Are Unidentified}}
\addto\captionsenglish{\renewcommand{\figurename}{Exhibit}}
\interfootnotelinepenalty=10000
\newcommand{\exhibitblock}[3]{\par\vspace{3pt}\noindent\begin{minipage}{\textwidth}\centering\includegraphics[width=0.94\textwidth]{#1}\captionof{figure}{\label{#2}#3}\end{minipage}\par\vspace{5pt}}
\title{Evaluating a Retraining Program\\When Its Losers Are Unidentified}
\author{Francesco Pinna\thanks{Department of Economics and Finance, LUISS Guido Carli. Email: \href{mailto:francesco.pinna@studenti.luiss.it}{francesco.pinna@studenti.luiss.it}. The policy, the pilot and all data are synthetic. The code and results that reproduce every number and exhibit are available at \href{https://github.com/francescopinna03/Unidentified-losers}{github.com/francescopinna03/Unidentified-losers}. All remaining errors are my own.}}
\date{October 29, 2026}
\begin{document}
\maketitle
\vspace{-1.2em}
\begin{abstract}
\noindent A public employment service must decide whether to extend a two-year retraining program, paired with a flat allowance of €2,000 per assigned worker, on the evidence of a randomized pilot. The lottery identifies an average earnings gain of €1,529 and, at an illustrative program cost of €3,000 per person, a marginal value of public funds of 0.68. Since no worker is observed both with and without the program, the share of workers whose disposable income falls after the allowance is not identified and ranges from 2.6\% if the program barely changes the ranking of workers to 26\% if it reshuffles that ranking and 71\% without restrictions. The same problem affects welfare measures that weight workers by their counterfactual earnings, while welfare changes computed from the two marginal distributions alone remain identified, so that the marginal value of public funds, once weighted toward the bottom of the counterfactual distribution, ranges from 0.66 to 1.30 across the matchings the data allow. Besides reporting these quantities under declared scenarios together with the divergence at which each claim breaks, the evaluation compares the allowance with a top-up of equal outlay that caps the share of losers under any matching at the price of weaker incentives to earn.\par
\vspace{4pt}\noindent\textit{Keywords}. Program evaluation, active labor market policy, distributional effects, partial identification, marginal value of public funds.\par
\noindent\textit{JEL codes}. D61, H43, J68.
\end{abstract}

\section{The policy question and the conceptual framework}
Retraining workers with low or unstable incomes is among the most widely funded active labor market policies in Europe, with a national program of this kind financed by the Italian recovery plan (PNRR), so whether such programs raise the earnings of the workers they reach, and at what fiscal cost, is a live question. The policy evaluated here is a two-year retraining program offered by a public employment service to workers with low or unstable incomes, with a flat allowance of €2,000 over the two years paid to those who attend. The program was piloted through a lottery of the offer among 6,000 eligible applicants, stratified by contract type, whose results form the evidence on which the decision to extend it rests. The paper asks three questions of that evidence, namely whether the program raises the earnings of the workers it reaches, what it delivers per euro of public money, and how many workers it leaves below the earnings they would have had without it. As Fernández and Rodrik \cite{fernandezrodrik} show, uncertainty about who would lose from a reform can block reforms that most would support ex post, a problem that the compensation principle studied by Schulz, Tsyvinski and Werquin \cite{schulz} can address only if losers can be found, so the share of losers is treated here as an evaluative object on the same footing as the average effect.

The literature gives reasons for both signs, since training can raise earnings through transferable skills, better matching and shorter unemployment spells, while Card, Kluve and Weber \cite{card} document that such gains appear with a lag and depend on the horizon and on the target population, whereas in the short run training programs show lock-in effects, since participants search less and work fewer hours during the course, to which a mismatch between the skills taught and local demand can add a lasting loss of earnings for some participants. Over a 24-month window that partly overlaps with the course, the forgone earnings of a worker who declines job offers to attend can therefore exceed the allowance, a heterogeneity that average effects can hide, as Bitler, Gelbach and Hoynes \cite{bitler} show. The framework is thus a trade-off between the average earnings gain, its fiscal cost and the losses of a minority, in which the open question for the evidence is which of the two forces dominates and for whom.

\section{Data and identification strategy}
The public employment service records assignment to the offer and attendance because it administers the lottery and pays the allowance, while earnings come from the monthly contribution records that employers file for each employee, which report pay and contract type because contributions are assessed on them, with tax returns for self-employment income and the unemployment insurance registers for benefit spells. The outcome is pre-tax labor earnings cumulated over the 24 months after assignment, excluding the allowance, so that the fiscal calculation can apply a tax rate to it, with a worker counted as a loser when disposable income, that is after-tax earnings plus the tax-free allowance, falls below the after-tax earnings she would have had without the program. Besides missing informal work, a limit that returns in Section 4 when income-contingent transfers raise the return to underreporting, these records do not report hours, so a fall in earnings cannot be separated into fewer days worked and lower pay.

The lottery makes the workers assigned to the offer comparable with the applicants who were not drawn, so the difference in cumulated earnings identifies the effect of the offer. In the synthetic pilot every assigned worker attends, so that the effect of the offer coincides with the effect of participation and both the program cost of €3,000 and the allowance are paid per assigned person. The estimated effect is €1,529, with an approximate 95\% interval from €1,270 to €1,789 against a population value of €1,575. The comparison can mislead in two ways, since retrained workers may find jobs at the expense of non-assigned applicants in the same local market, which Crépon, Duflo, Gurgand, Rathelot and Zamora \cite{crepon} document for job search assistance, so the pilot can overstate the gain that a scaled-up program would deliver. Workers who leave the registers, by moving abroad or into informal work, appear with zero recorded earnings, which biases the comparison if it happens more often in one arm. Exhibit~\ref{fig:pilot} reports the estimates by cell for stable part-time employees, fixed-term workers and intermittent workers, with untreated mean earnings of €24,000, €18,000 and €12,000 over two years and dispersion growing in that order.

\exhibitblock{figures/exhibit_1.pdf}{fig:pilot}{\textbf{Average effect of the offer and share of losers consistent with the pilot.} Panel A reports average effects by cell with 95\% Wald intervals. Panel B reports, with the two scenarios of Section 3, the identification bounds on the share of losers implied by the benchmark earnings distributions, which a larger pilot would estimate more precisely without narrowing, since no sample identifies the dependence between the two potential outcomes.}

\section{Costs, benefits and distributional effects}
Costs and benefits are summarized with the marginal value of public funds proposed by Hendren and Sprung-Keyser \cite{hendren}, the ratio between the beneficiaries' willingness to pay for the policy and its net cost to the government. With a program cost of €3,000 per assigned person, a tax rate of 30\% on earnings gains and the allowance valued at face value, willingness to pay is €3,071 per person against a net cost of €4,541, an MVPF of 0.68 with a sampling interval from 0.63 to 0.73 (Exhibit~\ref{tab:eval}). Without the allowance the ratio falls to 0.42, because the allowance is a pure transfer whose own MVPF is one. None of these numbers involves the distribution of gains and losses across individuals, which is where the information in the pilot ends.

The pilot observes different people in the two arms, so it identifies the distribution of earnings with the program and without it, never both for the same person. A matching pairs what each worker would earn with the program with what that same worker would earn without it. The share of workers who end up below the earnings they would have had, even after the allowance, depends on how the two are matched, a matching that no sample size can settle, as Fan and Park \cite{fanpark} show. The evaluation therefore states two reference scenarios before looking at the results. In the persistent scenario the course leaves the ranking of workers almost intact, so that whoever would have earned more without it also earns more with it except for 10\% of workers reshuffled at random, while in the independent scenario it reshuffles the ranking entirely. Under the persistent scenario 2.6\% of assigned workers lose despite the allowance, under the independent one 25.9\%, and with no restriction on the matching the share can reach 71.3\%, as Panel B of Exhibit~\ref{fig:pilot} shows cell by cell. Each scenario is then allowed to be wrong by a measured amount, the divergence of the truth from it, reported as a share of the divergence of the independent scenario from the persistent one.\footnote{Formally, the divergence between two matchings of the same earnings distributions is their Kullback-Leibler divergence, computed on a 25-point grid for each cell, on which $D_{\mathrm{KL}}(R_{\mathrm{ind}}\,\|\,R_{\mathrm{pers}})=2.09$ and the unrestricted share of losers is 67.6\%. The divergence is asymmetric, so the normalization is a reading aid and never a linear interpolation between the two scenarios. Shares of losers and shortfalls at the scenarios and without restrictions are continuous closed forms, while values at a positive divergence and the weighted measures are computed on the grid. On the computational side, the worst case within a given distance is a convex program, solved with the entropic transport methods of Léonard \cite{leonard} and of Galichon and Henry \cite{galichon}. Fan, Park and Xu \cite{fanparkxu} study the related problem of distributional model risk in marginal problems, with the marginals themselves in Wasserstein balls.} Exhibit~\ref{fig:losers} traces the worst-case share as that divergence grows. Under the persistent scenario the claim that at most 10\% lose survives as long as the divergence of the truth from that scenario stays below 5.8\% of the gap, whereas under the independent scenario it already fails at the scenario itself and only the weaker claim that at most half lose survives, up to 22\% of the same gap, in the spirit of the breakdown frontiers of Masten and Poirier \cite{masten}. With 5\% or 20\% reshuffled instead of 10\%, the share is 1.3\% or 5.2\% and the claim survives to 8.0\% or 2.7\% of the gap.

\exhibitblock{figures/exhibit_2_losers.pdf}{fig:losers}{\textbf{Worst-case share of losers as a function of the divergence from the reference scenario.} Panel A reports the worst-case share of losers as the truth diverges from each scenario, in percent of the gap between the two scenarios. Panel B reports the share of losers in each cell under the aggregate worst case at 5\% of that gap.}

In the aggregate worst case at 5\% of the gap the shares by cell are 6.1\%, 10.1\% and 12.9\% under the persistent scenario and 27.9\%, 40.0\% and 46.7\% under the independent one, so the intermittent workers, with the lowest earnings and the widest dispersion, are the most exposed under both. Applied to the severity of losses, the same computation shows that at either scenario a loser falls short of her counterfactual disposable income by about €3,280 on average, while the population average shortfall rises from €85 under the persistent scenario to €396 at 5\% of the gap, from €852 to €1,310 under the independent one, and to €1,671 without restrictions, against an average after-tax gain of €1,103.

The marginal value of public funds is identified because it aggregates willingness to pay without weights, whereas the distributional version used here weights each person by her counterfactual earnings, so that it averages the product of a weight that depends on counterfactual earnings and a gain that depends on realized earnings, which depends on how the two are matched. With weights that are twice the average at the bottom of the counterfactual distribution and zero at the top, the weighted MVPF is 0.69 under the persistent scenario and 0.98 under the independent one, within a range from 0.66 to 1.30 across all matchings that straddles one, while the unweighted measure stays at 0.68 without ambiguity. At 5\% of the gap the range is from 0.66 to 0.80 around the persistent scenario and from 0.83 to 1.13 around the independent one. With weights three times the average at the bottom, the full range widens to 0.65 to 1.67.

\section{Trade-offs and policy conclusion}
Raising every treated income by the same amount, the flat allowance protects nobody in particular and leaves its distributional effect entirely to the unobserved matching. The same expected outlay spent as a top-up indexed to after-tax earnings, which raises each worker's disposable income to the 67th percentile of counterfactual after-tax earnings in the cell, caps the share of losers at 33.2\% in every cell under any matching and cuts the worst expected shortfall from €1,671 to €660, because only workers whose counterfactual earnings exceed the threshold can lose, so that the same money buys a guarantee that holds under any matching of earnings distributions that do not respond to the transfer, in place of one that holds only under a declared scenario. The price is an implicit marginal tax rate of 100\% below the threshold, which weakens the incentive to earn and raises the return to underreporting income, the trade-off that Saez \cite{saez} analyzes for the optimal phase-out of income transfers. Since the top-up is indexed to after-tax earnings, every euro of gross earnings forgone by its recipients, 56\% of assigned workers, raises its net cost by one euro through a larger transfer and lower taxes while leaving willingness to pay unchanged to first order, so its MVPF of 0.68 falls by about 0.015 per €100 forgone per assigned worker and to 0.56 if recipients earned 10\% less. Verlaat, Todeschini and Ramos \cite{verlaat} find in the randomized evaluation of a guaranteed income in Barcelona that beneficiaries were 22\% less likely to work after two years and that a slower phase-out attenuates the effect, while Hyman, Kovak and Leive \cite{hyman} find that a wage insurance paid only to re-employed workers raises employment, so the direction of the behavioral response depends on whether the transfer is conditioned on working, which leaves the comparison between an assumption the data cannot test and a behavioral response that only a second randomized arm could measure.

\par\vspace{5pt}\noindent\begin{minipage}{\textwidth}\centering\footnotesize
\setlength{\tabcolsep}{3pt}\begin{tabular}{lrrr}
\toprule
& Program only & With allowance & Top-up, equal outlay \\
\midrule
Average earnings gain, € & 1,529 & 1,529 & 1,529 \\
MVPF, program cost 2,000 / 3,000 / 4,000 & 0.69 / 0.42 / 0.30 & 0.87 / 0.68 / 0.55 & 0.87 / 0.68 / 0.55 \\
Distributional MVPF, persistent / independent & 0.45 / 0.97 & 0.69 / 0.98 & matching-dependent \\
Distributional MVPF, all matchings & 0.39 to 1.54 & 0.66 to 1.30 & matching-dependent \\
Losers, persistent scenario & 4.0\% & 2.6\% & 2.4\% \\
Losers, independent scenario & 40.1\% & 25.9\% & 24.2\% \\
Losers, no restriction on the matching & 89.4\% & 71.3\% & at most 33.2\% \\
Expected shortfall, no restriction, € per person & 2,474 & 1,671 & 660 \\
\bottomrule
\end{tabular}
\captionof{figure}{\label{tab:eval}\textbf{Static policy comparison.} Synthetic benchmark, distributional MVPFs on the 25-point grid, amounts in euros over two years, earnings held fixed in every column, so that only the second column rests on the pilot and the other two are simulations. Losers and shortfalls are measured on disposable income. The MVPF uses a 30\% tax rate on earnings gains at program cost €3,000 unless stated and values transfers at face value, so the top-up column reports its MVPF before behavioral responses. The top-up is indexed to after-tax earnings with a threshold at the 67th percentile of counterfactual after-tax earnings in each cell and costs €2,000 per person in expectation.}
\end{minipage}\par\vspace{5pt}

The evidence supports with confidence a positive average earnings gain and an MVPF below one at a program cost of €3,000, a ratio that, since the allowance enters benefits and net costs equally, reaches one only at a program cost equal to the average gain, €1,529. Because the claim that few workers lose holds only under the persistent scenario and close to it, neither that claim nor the verdict of a welfare measure weighted by counterfactual earnings is a finding of the pilot, so an evaluation should publish both under at least two declared scenarios with the divergence at which each claim breaks, together with the incidence on the most vulnerable cell. Three limitations are specific to this design, the first being that the synthetic earnings are uniform within cells, without zero earners or a long right tail, two features that could materially change the thresholds and the shortfalls. The 24-month horizon lies partly inside the lock-in period, so some losers at 24 months may be winners at 48, which a longer follow-up in the same registers can establish. Displacement is invisible in a lottery among individuals and calls for randomizing the intensity of the offer across local markets before a roll-out. Where protecting losers is part of the purpose of the policy, the top-up at equal outlay makes explicit a trade-off between robustness to the unknown matching and incentives to earn that belongs in the decision together with the average gain.
\newpage
\begin{thebibliography}{99}
\small
\bibitem{bitler} Bitler, M. P., J. B. Gelbach and H. W. Hoynes (2006). What Mean Impacts Miss: Distributional Effects of Welfare Reform Experiments. \emph{American Economic Review}, 96(4), 988-1012. \href{https://doi.org/10.1257/aer.96.4.988}{doi:10.1257/aer.96.4.988}.
\bibitem{card} Card, D., J. Kluve and A. Weber (2018). What Works? A Meta Analysis of Recent Active Labor Market Program Evaluations. \emph{Journal of the European Economic Association}, 16(3), 894-931. \href{https://doi.org/10.1093/jeea/jvx028}{doi:10.1093/jeea/jvx028}.
\bibitem{crepon} Crépon, B., E. Duflo, M. Gurgand, R. Rathelot and P. Zamora (2013). Do Labor Market Policies Have Displacement Effects? Evidence from a Clustered Randomized Experiment. \emph{Quarterly Journal of Economics}, 128(2), 531-580. \href{https://doi.org/10.1093/qje/qjt001}{doi:10.1093/qje/qjt001}.
\bibitem{fanpark} Fan, Y. and S. S. Park (2010). Sharp Bounds on the Distribution of Treatment Effects and Their Statistical Inference. \emph{Econometric Theory}, 26(3), 931-951. \href{https://doi.org/10.1017/S0266466609990168}{doi:10.1017/S0266466609990168}.
\bibitem{fanparkxu} Fan, Y., H. Park and G. Xu (2025). Quantifying Distributional Model Risk in Marginal Problems via Optimal Transport. \emph{Mathematics of Operations Research}. \href{https://doi.org/10.1287/moor.2024.0557}{doi:10.1287/moor.2024.0557}.
\bibitem{fernandezrodrik} Fernández, R. and D. Rodrik (1991). Resistance to Reform: Status Quo Bias in the Presence of Individual-Specific Uncertainty. \emph{American Economic Review}, 81(5), 1146-1155.
\bibitem{galichon} Galichon, A. and M. Henry (2026). \emph{An econometrician's guide to optimal transport}. arXiv:2604.04227, version 2. \url{https://arxiv.org/abs/2604.04227}.
\bibitem{hendren} Hendren, N. and B. Sprung-Keyser (2020). A Unified Welfare Analysis of Government Policies. \emph{Quarterly Journal of Economics}, 135(3), 1209-1318. \href{https://doi.org/10.1093/qje/qjaa006}{doi:10.1093/qje/qjaa006}.
\bibitem{hyman} Hyman, B. G., B. K. Kovak and A. Leive (2024). \emph{Wage Insurance for Displaced Workers}. NBER Working Paper 32464. \href{https://doi.org/10.3386/w32464}{doi:10.3386/w32464}.
\bibitem{leonard} Léonard, C. (2014). A survey of the Schrödinger problem and some of its connections with optimal transport. \emph{Discrete and Continuous Dynamical Systems}, 34(4), 1533-1574. \href{https://doi.org/10.3934/dcds.2014.34.1533}{doi:10.3934/dcds.2014.34.1533}.
\bibitem{masten} Masten, M. A. and A. Poirier (2020). Inference on Breakdown Frontiers. \emph{Quantitative Economics}, 11(1), 41-111. \href{https://doi.org/10.3982/QE1288}{doi:10.3982/QE1288}.
\bibitem{saez} Saez, E. (2002). Optimal Income Transfer Programs: Intensive versus Extensive Labor Supply Responses. \emph{Quarterly Journal of Economics}, 117(3), 1039-1073. \href{https://doi.org/10.1162/003355302760193959}{doi:10.1162/003355302760193959}.
\bibitem{schulz} Schulz, K., A. Tsyvinski and N. Werquin (2023). Generalized Compensation Principle. \emph{Theoretical Economics}, 18(4), 1665-1710. \href{https://doi.org/10.3982/TE3971}{doi:10.3982/TE3971}.
\bibitem{verlaat} Verlaat, T., F. Todeschini and X. Ramos (2025). The employment effects of a means-tested guaranteed income policy. \emph{Journal of Public Economics}, 105420. \href{https://doi.org/10.1016/j.jpubeco.2025.105420}{doi:10.1016/j.jpubeco.2025.105420}.
\end{thebibliography}
\end{document}
